% !TEX root = ../tcc.tex

\begin{resumo}[Abstract]
 \begin{otherlanguage*}{english}
The CEDIS Ecosystem at the University of Brasília integrates multiple software applications to support teaching and learning, yet its analytics remain restricted to descriptive usage indicators and raw response counts without diagnosing students' conceptual mastery. This work aims to develop Cognita, a module integrated into the ecosystem that applies Bayesian Knowledge Tracing (BKT) to multi-source response data to support a dual feedback loop: indicating priority concepts for student review and signaling low-absorption concepts to instructors. Characterized as applied, mixed-methods (qualitative-quantitative), and exploratory research, the study is structured into two phases. Phase~1 (Undergraduate Thesis~I) encompassed a narrative literature review, technical source-code analysis of CEDIS Deck Manager, requirements elicitation and prioritization using an adapted Lean Inception and the URPS+ model, integration architecture design, logical relational data modeling compliant with Brazil's General Data Protection Law (LGPD), and high-fidelity prototyping. As outcomes of this phase, the study produced a characterization of the ecosystem's data and API contracts, a Minimum Viable Product (MVP) comprising 20 functional requirements and 17 non-functional requirements, an architecture and data model enforcing explicit consent and cascading data deletion, ten validated interface prototypes, and the project schedule. The findings indicate that the ecosystem provides the technical foundation needed for interpretable, LGPD-compliant knowledge tracing, consolidating the groundwork for Phase~2 (Undergraduate Thesis~II), which will cover ethics committee submission, BKT pipeline implementation, predictive evaluation against a baseline, instructor validation, and usability assessment.

\vspace{\onelineskip}

\noindent
\textbf{Keywords}: Knowledge Tracing. Bayesian Knowledge Tracing. Learning Analytics. CEDIS Ecosystem. Software Engineering.
 \end{otherlanguage*}
\end{resumo}
